\section*{Bab \ref{ch:b3:ode} --- Persamaan Diferensial
Biasa}

\begin{solution}{exo:b3:ode:1}
(a) Dengan memisahkan variabelnya: $x(t) = \frac1{1 - t}$ pada
$\intoo{-\infty}1$: yakni peledakan di $T_+ = 1$, dengan $x(t) \to
+\infty$ --- tepat \cref{thm:b3:ode:maximal}(2).
(b) $x(t) = \tan t$ pada $\intoo{-\pi/2}{\pi/2}$: yakni peledakan di
kedua ujungnya.
(c) $x(t) = \frac{1}{1 + \eu^{-t}}$, yang global: sebab solusinya
tinggal di $\intoo01$, yakni himpunan terbatas, sehingga grafiknya tak dapat
lolos dari setiap \omterm{def:b3:topology:compact}{kompak} $\R\times\R$ dalam waktu berhingga --- jadi $T_\pm
= \pm\infty$. Perhatikan bahwa (a) dan (b) tak bertentangan dengan
\cref{cor:b3:ode:globallinear}: sebab $x^2$ dan $1 + x^2$ bertumbuh
superlinear.
\end{solution}

\begin{solution}{exo:b3:ode:2}
(a) Inilah taksiran \cref{cor:b3:ode:uniqueness} dengan $t_0
= 0$. Untuk ketajamannya: bagi $x' = Lx$ (yang Lipschitz-$L$ secara global), dua
solusinya berselisih tepat $(x_0 - y_0)\eu^{Lt}$.
(b) Taksirannya berbunyi: pemetaan \omterm{ex:b3:ode:planeclassification}{aliran} pada waktu $t$ bersifat
Lipschitz-$\eu^{L\abs t}$ terhadap syarat awalnya ---
yakni \omterm{def:b3:topology:continuity}{kekontinuannya}, secara seragam untuk $t$ pada \omterm{def:b3:topology:compact}{kompak}; sedangkan pada himpunan terbatas
bagi $f$ yang tak Lipschitz secara global, jalankanlah hal yang sama pada sebuah tabung \omterm{def:b3:topology:compact}{kompak}
di sekitar \omterm{def:b3:topology:pathconnected}{lintasannya} dengan konstanta lokalnya.
\end{solution}

\begin{solution}{exo:b3:ode:3}
(a) Berlaku $\abs{\sin(tx)} \leq 1$: jadi terbatas, yakni pertumbuhan linear dengan
$\alpha = 0$ dan $\beta = 1$: sehingga \cref{cor:b3:ode:globallinear} pada
$U = \R\times\R$.
(b) Berlaku $\abs{tx/(1 + x^2)} \leq \abs t\cdot\frac12$: jadi kembali
sublinear (bahkan terbatas pada \omterm{def:b3:topology:compact}{kompak} waktunya): jadi global.
(c) Dengan $X = (x, x')$: $X' = \bigl(\begin{smallmatrix}0 & 1\\ -q(t)
& 0\end{smallmatrix}\bigr)X$: yakni linear dengan koefisien yang
\omterm{def:b3:topology:continuity}{kontinu}: jadi global (yakni latar
\cref{thm:b3:ode:linearstructure}).
(d) Global; dan lewat Grönwall seperti pada \cref{cor:b3:ode:globallinear}:
$\norm{x(t)} \leq \norm{x(t_0)}\,\eu^{M\abs{t - t_0}}$ dengan
$M = \sup\vertiii{A}$.
\end{solution}

\begin{solution}{exo:b3:ode:4}
(a) Berlaku $A^2 = -I$ untuk yang pertama: sehingga $\eu^{tA} = \cos t\,I + \sin
t\,A = \bigl(\begin{smallmatrix}\cos t & -\sin t\\ \sin t &
\cos t\end{smallmatrix}\bigr)$. Untuk blok Jordannya: $\lambda I$ dan
$N$ komutatif: sehingga $\eu^{tA} =
\eu^{\lambda t}\bigl(\begin{smallmatrix}1 & t\\ 0 &
1\end{smallmatrix}\bigr)$. Untuk yang ketiga: $A^2 = I$: sehingga $\eu^{tA} =
\cosh t\,I + \sinh t\,A$.

(b) Sistemnya $X' = \bigl(\begin{smallmatrix}0&1\\-1&0
\end{smallmatrix}\bigr)X + \bigl(\begin{smallmatrix}0\\
\cos\omega t\end{smallmatrix}\bigr)$; lalu Duhamel dengan
\omterm{thm:b3:ode:linearstructure}{resolven} rotasinya memberikan solusi khususnya: untuk
$\omega \neq 1$, $x_p(t) = \frac{\cos(\omega t)}{1 -
\omega^2}$ (periksalah langsung); sedangkan untuk $\omega = 1$ integral
$\int_0^t\sin(t - s)\cos s\,\dd s = \frac t2\sin t$ menghasilkan
pertumbuhan \emph{sekular} $x_p = \frac t2\sin t$: yakni resonansinya
--- sebab penggeraknya memompakan energi pada frekuensi alaminya dan
amplitudonya tumbuh secara linear.
\end{solution}

\begin{solution}{exo:b3:ode:5}
(a) Berlaku $w' = x_1x_2'' - x_1''x_2 = x_1(-px_2' - qx_2) - (-px_1'
- qx_1)x_2 = -p\,w$: sehingga $w(t) = w(t_0)\eu^{-\int_{t_0}^tp}$,
yang tak pernah nol atau identik nol. Lalu jika $w \neq 0$, vektor
$(x_i, x_i')(t_0)$ bebas di $\R^2$, dan karena
ruang solusinya berdimensi $2$
(\cref{thm:b3:ode:linearstructure} bagi sistemnya), maka $(x_1,
x_2)$ merupakan basis.

(b) Dengan $u = \int\frac{\eu^{-\int p}}{x_1^2}$: $x_2 = x_1u$,
$x_2' = x_1'u + \frac{\eu^{-\int p}}{x_1}$, dan
\[
x_2'' + px_2' + qx_2
= u\,(x_1'' + px_1' + qx_1) +
\Bigl(-\,p\frac{\eu^{-\int p}}{x_1} +
p\frac{\eu^{-\int p}}{x_1}\Bigr) = 0
\]
(sebab suku silangnya meniadakan diri secara persis; uraikanlah dengan saksama). Untuk
$t^2x'' - 2x = 0$, yakni $x'' - \frac2{t^2}x = 0$ (dengan $p = 0$)
dan $x_1 = t^2$: $u = \int t^{-4} = -\frac1{3t^3}$, sehingga $x_2
= -\frac1{3t}$: jadi solusi umumnya $at^2 + \frac bt$.
\end{solution}

\begin{solution}{exo:b3:ode:6}
Kesetimbangannya $0, 1$; dan $f'(x) = 1 - 2x$: sehingga $f'(0) = 1 > 0$
(jadi tak \omterm{def:b3:ode:stability}{stabil} --- sebab solusi di dekatnya menjauh, seperti ditunjukkan
bentuk eksplisitnya), sedangkan $f'(1) = -1 < 0$: jadi \omterm{def:b3:ode:stability}{stabil} asimtotik
(\cref{thm:b3:ode:linearization} dalam dimensi $1$). Untuk
$x(0) \in \intoo01$: $f > 0$ di sana, sehingga selama
solusinya tinggal di $\intoo01$ ia naik; dan ia tak pernah dapat mencapai
$0$ atau $1$ (menurut ketunggalannya: sebab keduanya \omterm{def:b3:topology:pathconnected}{lintasan}), sehingga ia tinggal,
terbatas --- jadi global --- dan naik ke sebuah limit $L
\in \intoc{x(0)}1$. Seandainya $f(L) \neq 0$, maka $x' \geq c > 0$
di dekat limitnya, yang memaksa $x$ melewati $L$: jadi $f(L) = 0$ dan $L = 1$;
sedangkan setangkup dengan itu $x \to 0$ di $-\infty$. Dan secara eksplisit $x(t) =
\frac1{1 + C\eu^{-t}}$ membenarkan segalanya. Secara umum: jika sebuah
solusi otonom skalar mempunyai $x'(t_0) = 0$, maka $x(t_0)$
merupakan kesetimbangan dan ketunggalannya membuat $x$ konstan;
sedangkan selain itu $f(x(t))$ menjaga tanda yang tetap (sebab ia tak pernah lenyap,
dan $t \mapsto f(x(t))$ \omterm{def:b3:topology:continuity}{kontinu}): jadi $x$ monoton
sejati --- sehingga solusi periodik yang tak konstan bersifat
mustahil.
\end{solution}

\begin{solution}{exo:b3:ode:7}
(a) $\frac{\dd}{\dd t}H(x,y) = H_xx' + H_yy' = H_xH_y +
H_y(-H_x) = 0$.
(b) Himpunan aras $H = \frac{y^2}2 + \frac{x^4}4$ bersifat
\omterm{def:b3:topology:compact}{kompak} (sebab $H$ koersif), sehingga solusinya terjebak di \omterm{def:b3:topology:compact}{kompak}:
jadi global dan terbatas (\cref{thm:b3:ode:maximal}). Untuk kestabilan
$(0,0)$: $V = H$ bersifat definit positif (sebab $H = 0$ hanya di
titik asalnya) dengan $\dot V = 0$: sehingga \cref{thm:b3:ode:lyapunov}. Sedangkan
pelinearannya $x' = y$, $y' = 0$ mempunyai matriks nilpoten
yang tak terdiagonalkan bernilai eigen $0$:
sehingga \cref{thm:b3:ode:linearization} bungkam (sebab hipotesisnya
$\operatorname{Re}\lambda < 0$ gagal), dan sungguh
sistem terlinearkannya tak \omterm{def:b3:ode:stability}{stabil} (sebab $y_0 \ne 0$ menghanyut) sedangkan
yang taklinear \omterm{def:b3:ode:stability}{stabil}: jadi pelinearan pada kesetimbangan yang tak
hiperbolik tak membuktikan apa pun.
\end{solution}

\begin{solution}{exo:b3:ode:8}
(a) Berlaku $\dot V = y\,y' + \sin x\cdot x' = y(-\sin x - cy) +
y\sin x = -cy^2 \leq 0$, dan $V = \frac{y^2}2 + (1 - \cos x)$
bersifat definit positif pada $\{\abs x < 2\pi\}$ di sekitar
titik asalnya: jadi \omterm{def:b3:ode:stability}{stabil} (\cref{thm:b3:ode:lyapunov}).
(b) Matriks terlinearkannya di $(0,0)$ adalah
$\bigl(\begin{smallmatrix}0 & 1\\ -1 &
-c\end{smallmatrix}\bigr)$, dengan polinomial karakteristik
$\lambda^2 + c\lambda + 1$: yang berakar $\frac{-c \pm \sqrt{c^2 -
4}}2$ --- keduanya real negatif jika $c \geq 2$, dan kompleks berbagian
real $-\frac c2 < 0$ jika $0 < c < 2$. Jadi dalam segala kasusnya
$\operatorname{Re}\lambda < 0$: sehingga
\cref{thm:b3:ode:linearization} memberikan kestabilan asimtotiknya
(walaupun $\dot V$-nya merosot).
(c) Di $(\pi, 0)$: $\sin(\pi + u) = -\sin u$, dengan pelinearannya
$\bigl(\begin{smallmatrix}0&1\\1&-c\end{smallmatrix}\bigr)$,
dan karakteristiknya $\lambda^2 + c\lambda - 1$: yang berakar bertanda
berlawanan (sebab $\lambda_+\lambda_- = -1$). Lalu sepanjang vektor eigen
tak \omterm{def:b3:ode:stability}{stabilnya}, sistem linearnya mempunyai solusi yang lolos secara
eksplisit $\eu^{\lambda_+t}v_+$ dengan $\lambda_+ > 0$: sehingga
bandul terbalik itu tak \omterm{def:b3:ode:stability}{stabil} untuk setiap redamannya.
\end{solution}

\begin{solution}{exo:b3:ode:9}
Untuk $\alpha \in \intoo01$: selain $x \equiv 0$, setiap
\[
x_c(t) = \begin{cases}0 & t \leq c,\\
\bigl((1-\alpha)(t - c)\bigr)^{1/(1-\alpha)} & t \geq c,
\end{cases}
\]
bersifat $\mathcal C^1$ dan menyelesaikan persamaannya (sebab eksponennya
$\frac1{1-\alpha} > 1$ membuat turunannya lenyap di $c$): jadi sebuah
kontinum solusi yang melalui $(0,0)$. Sedangkan untuk $\alpha = 1$: $x
\mapsto \abs x$ bersifat Lipschitz-$1$ secara global
(sebab $\abs{\abs a - \abs b} \leq \abs{a - b}$), sehingga
\cref{thm:b3:ode:cauchylipschitz} berlaku dan satu-satunya
solusi yang melalui $0$ adalah $x \equiv 0$. Jadi perbatasannya tepat
syarat \omterm{def:b3:ode:lipschitz}{Lipschitz lokal} di $0$: sebab $\abs x^\alpha$ berhasil bagi
selisih yang tak terbatas di sana untuk $\alpha < 1$.
\end{solution}

\begin{solution}{exo:b3:ode:10}
(a) Andaikan $\norm{x(t_2)} > R$ untuk suatu $t_2 > 0$ dengan
$\norm{x(0)} \leq R$, lalu misalkan $t_1 = \sup\{t \leq t_2 :
\norm{x(t)} \leq R\}$: maka $\norm{x(t_1)} = R$ dan
$\norm{x(t)} > R$ pada $\intoc{t_1}{t_2}$. Lalu pada selang itu
$g(t) = \norm{x(t)}^2$ mempunyai $g'(t) = 2\langle x(t),
F(x(t))\rangle \leq 0$ (sebab hipotesisnya berlaku: karena $\norm{x(t)}
\geq R$), sehingga $g(t_2) \leq g(t_1) = R^2$: jadi bertentangan. Sehingga
bolanya invarian secara positif.
(b) Sebuah solusi yang terjebak di bola \omterm{def:b3:topology:compact}{kompaknya} tak dapat mempunyai
$T_+ < \infty$ (\cref{thm:b3:ode:maximal}(2)): jadi global
maju. Untuk sistem gradiennya: $\frac{\dd}{\dd t}G(x(t)) =
\langle\nabla G, -\nabla G\rangle = -\norm{\nabla G(x(t))}^2
\leq 0$: sehingga $G$ menurun, jadi solusinya tinggal di $\{G \leq
G(x_0)\}$, yang terbatas (menurut kekoersifannya: sebab di luar sebuah bola
yang besar, $G > G(x_0)$) dan tertutup: jadi \omterm{def:b3:topology:compact}{kompak}. Sehingga pelolosannya
mustahil: jadi setiap solusi sistem gradien yang koersif bersifat
global maju, yang meluncur menurun selamanya.
\end{solution}

\begin{solution}{exo:b3:ode:11}
(a) Untuk lema pembandingannya: misalkan $w = x - y$ pada selang
bersamanya; maka $w(0) \geq 0$ dan $w' = x' - y' \geq F(x) - F(y) =
c(t)w$ dengan $c(t) = \frac{F(x) - F(y)}{x - y}$ yang terbatas pada
selang waktu \omterm{def:b3:topology:compact}{kompaknya} (sebab $F$ \omterm{def:b3:ode:lipschitz}{Lipschitz lokal}); lalu
$(w\eu^{-\int c})' \geq 0$, sehingga $w \geq 0$ di sepanjangnya. Dengan
$F(x) = x^2$: $y(t) = \frac1{1 - t}$ meledak di $1$, dan $x
\geq y$ selama $x$ hidup; sehingga jika $T_+ > 1$, maka $x$ akan
berhingga di $t = 1$ padahal mendominasi $y \to \infty$:
jadi mustahil. Sehingga $T_+ \leq 1$.

(b) Untuk pembandingan terbaliknya (lewat lema yang sama dengan peran yang ditukar):
pada $\intcc01\cap\intco0{T_+}$, $t^2 \leq 1$ memberikan $x' \leq
x^2 + 1$, sedangkan $z(t) = \tan(t + \frac\pi4)$ memenuhi $z' =
z^2 + 1$ dengan $z(0) = 1 = x(0)$: sehingga $x \leq z$ selama
keduanya terdefinisi. Lalu karena $z$ berhingga pada
$\intco0{\frac\pi4}$, maka $x$ tak dapat meledak sebelum
$\frac\pi4$: jadi $T_+ \geq \frac\pi4$.

(c) Bersama-sama: $\frac\pi4 \leq T_+ \leq 1$ (dan secara numerik $T_+
\approx 0.96$). Moralnya: bahwa untuk $x' = F(t, x)$ dengan $F$ yang
superlinear terhadap $x$, solusinya meledak dalam waktu berhingga
setiap kali $\int^\infty\frac{\dd s}{F(s)} < \infty$ (sebab
solusi pembandingnya mencapai tak hingga dalam waktu berhingga itu);
sedangkan pertumbuhan linear, yang integralnya divergen, memaksa keberadaan
globalnya (\cref{exo:b3:ode:3}). Jadi ia integral Osgood yang
sama seperti pada \cref{pb:b3:complete:1}, yang kini mengatur
pelolosan ke tak hingga alih-alih pelolosan dari nol.
\end{solution}

\begin{solution}{exo:b3:ode:12}
(a) $W' = uv'' - u''v = u(-q_2v) - (-q_1u)v = (q_1 -
q_2)\,uv$.

(b) Misalkan $a < b$ akar berurutan $u$; lalu normalkan $u >
0$ pada $\intoo ab$ (sehingga $u'(a) > 0$ dan $u'(b) < 0$ --- yang tak nol
menurut ketunggalannya, sebab $u(a) = u'(a) = 0$ akan memaksa $u
\equiv 0$). Andaikan $v$ tak berakar di $\intoo ab$;
lalu normalkan $v > 0$ di sana (sehingga $v(a), v(b) \geq 0$ menurut
\omterm{def:b3:topology:continuity}{kekontinuannya}). Lalu integralkan (a):
\[
W(b) - W(a) = \int_a^b(q_1 - q_2)\,uv \;\leq\; 0 .
\]
Padahal $W(a) = u(a)v'(a) - u'(a)v(a) = -u'(a)v(a) \leq 0$ dan
$W(b) = -u'(b)v(b) \geq 0$: sehingga $W(b) - W(a) \geq 0$.
Jadi kesamaannya di sepanjangnya: sebab $\int(q_1 - q_2)uv = 0$ dengan $uv > 0$
pada selang terbukanya memaksa $q_1 = q_2$ di sana; dan $W(a) =
W(b) = 0$ memaksa $v(a) = v(b) = 0$; lalu $W \equiv 0$ pada
$\intcc ab$ (sebab turunannya lenyap), yakni
$(v/u)' = -W/u^2 = 0$ pada $\intoo ab$: jadi $v \propto u$.

(c) Ambillah $q_1 = m$ dan $u = \sin(\sqrt m(t - t_0))$, yang
akar berurutannya berjarak $\pi/\sqrt m$, dan $q_2 = q \geq
m$: maka menurut (b), setiap solusi $v$ dari $v'' + qv = 0$ lenyap di
setiap selang terbuka berpanjang $\pi/\sqrt m$ (sedangkan pada
alternatif merosotnya $q \equiv m$, $v \propto u$ lenyap
pula). Sedangkan jika $q \leq 0$: terapkan (b) dengan $q_1 = q$, $u =
v$, dan $q_2 = 0$ beserta solusi bebas akar $\mathbf 1$ dari
$v'' = 0$. Seandainya $v$ berakar dua kali secara berurutan, (b) akan memaksa
entah sebuah akar $\mathbf 1$ di antaranya entah kasus
merosotnya $\mathbf 1 \propto v$ --- keduanya mustahil: jadi $v$ lenyap
paling banyak sekali. Untuk ujinya: bagi $q = 1$, $\sin t$ lenyap setiap $\pi
= \pi/\sqrt1$, seperti yang diramalkan; sedangkan bagi $q = -1$, $\sinh t$
lenyap tepat sekali dan $\eu^t$ tak pernah --- jadi paling banyak satu
akar, seperti yang diramalkan.
\end{solution}

\begin{solution}{pb:b3:ode:1}
\textbf{1.} Berlaku $\dot E = yy' + \sin x\cdot x' = -y\sin x +
y\sin x = 0$: sehingga energinya merupakan integral pertama. Lalu pada sebuah
\omterm{thm:b3:ode:maximal}{solusi maksimal}, $y^2 = 2(E_0 + \cos x) \leq 2(E_0 + 1)$: jadi $y$
terbatas; sehingga $\abs{x(t)} \leq \abs{x(0)} +
t\sup\abs y$ tumbuh paling banyak secara linear: jadi pada sembarang selang waktu
berhingga \omterm{def:b3:topology:pathconnected}{lintasannya} tinggal di sebuah \omterm{def:b3:topology:compact}{kompak} $\R^2$, sehingga
\cref{thm:b3:ode:maximal}(2) memaksa $T_\pm = \pm\infty$.
Kesetimbangannya $(k\pi, 0)$; dengan pelinearan
$\bigl(\begin{smallmatrix}0&1\\
\mp1&0\end{smallmatrix}\bigr)$ karena $-\cos(k\pi) = \mp1$:
sehingga nilai eigennya $\pm\iu$ bagi $k$ yang genap (yakni bertipe pusat, yang tak menyimpulkan
dengan sendirinya) dan $\pm1$ bagi $k$ yang gasal (yakni pelana).

\textbf{2.} Karena $E$ konstan sepanjang solusinya, setiap
\omterm{def:b3:topology:pathconnected}{lintasannya} terkurung di sebuah himpunan aras $\{y^2 = 2(E_0 + \cos x)\}$. Untuk
$E_0 = -1$: hanya titik $(2k\pi, 0)$. Untuk $-1 < E_0 <
1$: dengan menulis $E_0 = -\cos a$, himpunannya berupa gabungan saling lepas
kurva tertutup $y = \pm\sqrt{2(\cos x - \cos a)}$ atas $x \in
[2k\pi - a, 2k\pi + a]$, satu di sekitar setiap kesetimbangan \omterm{def:b3:ode:stability}{stabilnya}.
Untuk $E_0 = 1$: kurva $y = \pm2\cos\frac x2$ yang menghubungkan
pelana berurutan --- yakni separatriksnya --- beserta
pelananya sendiri. Sedangkan untuk $E_0 > 1$: dua grafik $y =
\pm\sqrt{2(E_0 + \cos x)}$, yang terdefinisi untuk setiap $x$ dan tak pernah
menyentuh $y = 0$.

\textbf{3.} Fungsi $V = E + 1 = \frac{y^2}2 + (1 - \cos x)$
lenyap di $(0,0)$, positif pada sebuah \omterm{def:b3:topology:topology}{persekitaran} tertusuk
(dengan $\abs x < 2\pi$), dan $\dot V = 0 \leq 0$:
sehingga \cref{thm:b3:ode:lyapunov} memberikan kestabilannya. Namun tak asimtotik:
sebab solusi yang melalui $(a, 0)$ (dengan $0 < a$ yang kecil) tetap pada
kurva aras $E = -\cos a$, yang jaraknya ke titik asalnya
positif (sebab kurvanya memotong sumbu-$x$ hanya di $\pm a$):
sehingga $\varphi_t(a, 0) \not\to (0,0)$.

\textbf{4.} Pada kurva aras $C_a$: tak ada kesetimbangan (sebab $y = 0$
memaksa $x = \pm a$ dengan $\sin(\pm a) \neq 0$ untuk $0 < a <
\pi$), sehingga lajunya $\norm{(y, -\sin x)}$ berminimum positif
$m$ pada $C_a$ yang \omterm{def:b3:topology:compact}{kompak}. Lalu ikutilah solusinya dari
$(a, 0)$: pada setengah bidang bawahnya $x' = y < 0$, sehingga $x$
turun dari $a$ ke $-a$ dalam waktu berhingga (sebab integral seperempat
atau setengah periodenya konvergen: lihat analisis pertanyaan 5), lalu mencapai
$(-a, 0)$; lalu menurut kesetangkupan $(x, y) \mapsto (x, -y)$ dan $t
\mapsto -t$ pada persamaannya, setengah atasnya dilalui
kembali dalam waktu $\frac T2$ yang sama: sehingga solusinya kembali ke
$(a, 0)$ pada waktu $T$. Lalu ketunggalannya
(\cref{cor:b3:ode:uniqueness}) merambatkannya: $x(t + T) =
x(t)$ untuk setiap $t$: jadi periodik.

\textbf{5.} Pada cabang yang $y > 0$: $y = \sqrt{2(\cos x
- \cos a)}$ dan $\dd t = \frac{\dd x}{y}$; lalu mengintegralkan $x$
dari $-a$ ke $a$ memberikan setengah periodenya, dan kesetangkupan $x
\mapsto -x$ memparuhkan integralnya lagi:
\[
T(a) = 4\int_0^a\frac{\dd x}{\sqrt{2(\cos x - \cos a)}} .
\]
Untuk kekonvergenannya di $x = a^-$: $\cos x - \cos a = \sin(a)(a - x) +
O((a-x)^2)$ dengan $\sin a > 0$: sehingga integrannya berperilaku seperti
$\bigl(2\sin a\,(a - x)\bigr)^{-1/2}$, yang \omterm{def:b3:lebesgue:l1}{terintegralkan}.

\textbf{6.} Dengan $\sin\frac x2 = k\sin\varphi$ dan $k =
\sin\frac a2$: $\cos x - \cos a = 2(k^2 - \sin^2\frac x2)
= 2k^2\cos^2\varphi$ dan $\frac12\cos\frac x2\,\dd x =
k\cos\varphi\,\dd\varphi$, sehingga
\[
T(a) = 4\int_0^{\pi/2}
\frac{\dd\varphi}{\sqrt{1 - k^2\sin^2\varphi}} .
\]
Lalu saat $a \to 0^+$, $k \to 0$: sebab untuk $k \leq k_0 < 1$
integrannya terdominasi oleh $(1 - k_0^2\sin^2\varphi)^{-1/2}$,
yang \omterm{def:b3:topology:continuity}{kontinu} pada $\intcc0{\pi/2}$: sehingga kekonvergenan terdominasinya memberikan $T \to
4\cdot\frac\pi2 = 2\pi$. Lalu dengan menguraikan $(1 - u)^{-1/2} = 1 +
\frac u2 + \frac{3u^2}8 + \cdots$ dengan $u =
k^2\sin^2\varphi$ (lewat kekonvergenan normalnya untuk $k < 1$) dan
$\int_0^{\pi/2}\sin^2 = \frac\pi4$:
\[
T(a) = 2\pi\Bigl(1 + \frac{k^2}{4} + O(k^4)\Bigr) :
\]
sehingga keisokronannya berlaku hanya sampai orde pertamanya; dan periodenya tumbuh seiring
amplitudonya.

\textbf{7.} Saat $k \uparrow 1$ integrannya naik menuju $(1
- \sin^2\varphi)^{-1/2} = \frac1{\cos\varphi}$, yang
integralnya divergen: sehingga lewat kekonvergenan monoton, $T(a) \to
+\infty$ saat $a \to \pi^-$.

\textbf{8.} Pada cabang $y = 2\cos\frac x2$ (dengan $\abs x <
\pi$): $\frac{\dd x}{2\cos(x/2)} = \dd t$; lalu dengan $u = \frac
x2$, $\int\frac{\dd u}{\cos u} =
\ln\tan\bigl(\frac u2 + \frac\pi4\bigr)$, sehingga $t =
\ln\tan\bigl(\frac x4 + \frac\pi4\bigr)$, yakni
\[
x(t) = 4\arctan(\eu^t) - \pi,
\qquad y(t) = x'(t) = \frac{4\eu^t}{1 + \eu^{2t}} =
\frac{2}{\cosh t} .
\]
Untuk pemeriksaannya lewat energinya: dengan $\theta = \arctan\eu^t$,
$\sin2\theta = \frac1{\cosh t}$, sehingga $\cos x = -\cos4\theta =
-1 + \frac{2}{\cosh^2t}$ dan $E = \frac{y^2}2 - \cos x =
\frac2{\cosh^2t} + 1 - \frac2{\cosh^2t} = 1$: jadi \omterm{def:b3:topology:pathconnected}{lintasannya}
terletak pada separatriksnya, dan menurunkan $y^2 = 2(1 +
\cos x)$ di tempat $y > 0$ menghasilkan kembali $y' = -\sin x$. Limitnya: $x
\to \pm\pi$ dan $y \to 0$ saat $t \to \pm\infty$.

\textbf{9.} \omterm{def:b3:groups:action}{Orbitnya} menuju pelana $(\pi, 0)$ ke depan
dan $(-\pi, 0)$ ke belakang tetapi tak pernah tiba: sebab seandainya ia mencapai
$(\pi, 0)$ pada waktu berhingga $t^*$, maka dua \omterm{thm:b3:ode:maximal}{solusi maksimal} yang berbeda
--- yakni solusi separatriksnya dan solusi konstan
di pelananya --- akan melalui titik yang sama
$(t^*, (\pi, 0))$, yang bertentangan dengan
\cref{cor:b3:ode:uniqueness}. Jadi pelananya hanya didekati
secara asimtotik.

\textbf{10.} Untuk $E_0 > 1$: $y^2 = 2(E_0 + \cos x) \geq
2(E_0 - 1) > 0$: sehingga $y$ menjaga tandanya, dan $\abs{x'} = \abs y
\geq \sqrt{2(E_0 - 1)}$: jadi $x$ monoton sejati dan global
(pertanyaan 1), dengan $x(t) \to \pm\infty$. Lalu karena $y(t) =
\pm\sqrt{2(E_0 + \cos x(t))}$ dan $\cos$ berperiode $2\pi$,
maka $y$ kembali ke nilainya setiap kali $x$ maju sejauh $2\pi$;
dan waktu yang diperlukannya adalah
\[
\tau(E_0) = \int_{x_0}^{x_0 + 2\pi}\frac{\dd
x}{\sqrt{2(E_0 + \cos x)}} =
\int_{-\pi}^{\pi}\frac{\dd x}{\sqrt{2(E_0 + \cos x)}}
\]
(lewat substitusi dan keperiodikannya): sehingga $y$ berperiode $\tau$ ---
jadi bandulnya berpusar dengan laju rotasi yang asimtotik tetap
$2\pi/\tau \approx \sqrt{2E_0}$ untuk energi yang besar.

\textbf{11.} Metodenya, berurutan: \emph{energinya} (yakni $\dot E =
0$) menyusutkan \omterm{ex:b3:ode:planeclassification}{aliran} dua dimensinya menjadi kurva aras
satu dimensi; lalu \emph{keterbatasan} $y$ pada setiap arasnya ditambah
\omterm{thm:b3:ode:maximal}{pelolosan dari kompaknya} memberikan keberadaan globalnya; lalu
\emph{kekompakan} aras tertutupnya memberikan batas lajunya
sehingga keperiodikannya; lalu \emph{ketunggalannya} mengubah
kembalinya yang pertama menjadi keperiodikan yang persis, melarang kedatangan
pada pelananya dalam waktu berhingga, dan memisahkan tipe \omterm{def:b3:groups:action}{orbitnya};
lalu \emph{pelinearan dan Lyapunov} menggolongkan kesetimbangannya;
sedangkan \emph{integral periodenya} dianalisis lewat teorema kekonvergenan
\cref{ch:b3:lebesgue}. Jadi tak pernah kita memiliki --- atau
memerlukan --- solusi umum berbentuk tertutup: sebab
teori kualitatifnya memetik setiap gejala geraknya
dari persamaannya sendiri.

\textbf{12.} Integral parsial: $W_n = \int_0^{\pi/2}\sin^{2n-1}\varphi
\cdot\sin\varphi\,\dd\varphi = (2n-1)\int_0^{\pi/2}
\sin^{2n-2}\varphi\cos^2\varphi\,\dd\varphi = (2n-1)(W_{n-1}
- W_n)$, sehingga $W_n = \frac{2n-1}{2n}W_{n-1}$; lalu dengan $W_0 =
\frac\pi2$ dan $\prod_{j=1}^n\frac{2j-1}{2j} =
\frac{(2n)!}{4^n(n!)^2}$ (sebab $(2n)! = 2^nn!\prod(2j-1)$),
induksinya memberikan nilai yang ditampilkan itu. Untuk deret binomialnya: $(1 -
u)^{-1/2} = \sum_{n\geq0}c_nu^n$ dengan $c_n =
\frac{(2n)!}{4^n(n!)^2} \in \intoc01$ dan jari-jari $1$. Lalu untuk $k
\leq k_0 < 1$ dan $u = k^2\sin^2\varphi$, deret
$\sum c_nk^{2n}\sin^{2n}\varphi$ konvergen normal terhadap
$\varphi$ (sebab $c_nk^{2n} \leq k_0^{2n}$), sehingga pengintegralan suku demi
sukunya pada rumus pertanyaan 6 bersifat sah:
\[
T(a) = 4\sum_{n\geq0}c_nk^{2n}W_n
= 4\cdot\frac\pi2\sum_{n\geq0}c_n^2\,k^{2n}
= 2\pi\sum_{n\geq0}
\Bigl(\frac{(2n)!}{4^n(n!)^2}\Bigr)^{\!2}k^{2n} .
\]
Lalu dengan $c_0 = 1$, $c_1 = \frac12$, dan $c_2 = \frac38$: $T(a) =
2\pi\bigl(1 + \frac{k^2}4 + \frac{9k^4}{64} + O(k^6)\bigr)$,
dengan sisanya seragam untuk $k \leq k_0$ (sebab ekornya terdominasi sebuah
deret geometri).

\textbf{13.} Berlaku $k = \sin\frac a2 = \frac a2 - \frac{a^3}{48} +
O(a^5)$, sehingga
\[
k^2 = \frac{a^2}4 - \frac{a^4}{48} + O(a^6),
\qquad
k^4 = \frac{a^4}{16} + O(a^6) ,
\]
dan
\[
\frac{T(a)}{2\pi} = 1 + \frac14\Bigl(\frac{a^2}4 -
\frac{a^4}{48}\Bigr) + \frac9{64}\cdot\frac{a^4}{16} +
O(a^6)
= 1 + \frac{a^2}{16} + \frac{11\,a^4}{3072} + O(a^6) ,
\]
sebab $-\frac1{192} + \frac9{1024} = \frac{-16 + 27}{3072}
= \frac{11}{3072}$.

\textbf{14.} Pada bentuk eliptik pertanyaan 6, $a \mapsto
k = \sin\frac a2$ merupakan bijeksi \omterm{def:b3:topology:continuity}{kontinu} yang naik
sejati dari $\intoo0\pi$ ke $\intoo01$, dan untuk setiap
$\varphi \in \intoc0{\pi/2}$ integrannya $(1 -
k^2\sin^2\varphi)^{-1/2}$ naik sejati terhadap $k$: sehingga $T$
naik sejati. Untuk \omterm{def:b3:topology:continuity}{kekontinuannya}: pada $k \leq k_0 < 1$
integrannya terdominasi oleh $(1 -
k_0^2\sin^2\varphi)^{-1/2}$ yang \omterm{def:b3:topology:continuity}{kontinu}, sehingga kekonvergenan terdominasinya
berlaku sepanjang $k' \to k$. Lalu dengan limitnya $T \to 2\pi$ saat $a
\to 0^+$ (pertanyaan 6) dan $T \to +\infty$ saat $a \to \pi^-$
(pertanyaan 7), kemonotonan sejatinya dan teorema nilai
antaranya membuat $T$ menjadi bijeksi dari $\intoo0\pi$ ke
$\intoo{2\pi}{+\infty}$.

\textbf{15.} Sebuah jam mencacah ayunan; dan jika diatur pada amplitudo
yang melenyap, ia membukukan periode harmoniknya $2\pi$ per ayunan (dalam
satuan waktu bandulnya). Sedangkan jika dijalankan pada amplitudo $a$, periode
sebenarnya adalah $T(a) = 2\pi\bigl(1 + \frac{a^2}{16} +
O(a^4)\bigr)$: sehingga jamnya membukukan $2\pi$ padahal $T(a)$ sungguh
berlalu, jadi ia tertinggal sebesar pecahan
\[
\frac{T(a) - 2\pi}{T(a)} = \frac{a^2}{16} + O(a^4) .
\]
Untuk $a = 0.2$ rad (yakni sekitar $11.5$ derajat): $\frac{a^2}{16} =
\frac{0.04}{16} = \frac1{400}$, dan satu hari berdurasi $86400$ s: sehingga
jamnya kehilangan $86400/400 = 216$ detik --- yakni sekitar tiga setengah
menit --- per hari. Karena itu kedua obat historisnya:
memaksakan amplitudo tetap yang mungil (lewat rodagigi pelepasnya), atau membengkokkan
kendalanya sehingga periodenya sungguh bebas amplitudo
(lewat pipi sikloid Huygens, 1657).

\textbf{16.} Pada $\tau(E_0) = \int_{-\pi}^{\pi}\frac{\dd
x}{\sqrt{2(E_0 + \cos x)}}$ integrannya, untuk setiap
$x$ yang tetap, turun sejati terhadap $E_0$: sehingga $\tau$ turun
sejati. Lalu saat $E_0 \downarrow 1$ integrannya naik
secara titik demi titik menuju $\bigl(2(1 + \cos x)\bigr)^{-1/2} =
\frac1{2\abs{\cos\frac x2}}$, yang integralnya atas
$\intoo{-\pi}\pi$ divergen (sebab $\cos\frac x2$ lenyap berorde
pertama di $\pm\pi$): sehingga kekonvergenan monotonnya memberikan $\tau(E_0)
\to +\infty$. Sedangkan saat $E_0 \to +\infty$:
\[
\sqrt{2E_0}\,\tau(E_0) = \int_{-\pi}^{\pi}
\frac{\dd x}{\sqrt{1 + \cos x/E_0}} \longrightarrow 2\pi
\]
lewat kekonvergenan terdominasi (sebab untuk $E_0 \geq 2$ integrannya
paling banyak $\sqrt2$). Jadi $\tau \approx 2\pi/\sqrt{2E_0}$: sehingga satu
putaran memakan waktu rotasi bebas berlaju $\sqrt{2E_0}$,
dengan potensialnya menyusut menjadi riak --- yang cocok dengan laju rotasi
pertanyaan 10.

\textbf{17.} Sumbunya membawa solusi eksplisit $t
\mapsto (x_0\eu^{\alpha t}, 0)$ dan $t \mapsto (0,
y_0\eu^{-\gamma t})$, beserta kesetimbangan
$(0, 0)$: jadi keduanya gabungan \omterm{def:b3:groups:action}{orbit}. Sedangkan medannya bersifat $\mathcal
C^1$, jadi \omterm{def:b3:ode:lipschitz}{Lipschitz lokal}; sehingga solusi yang bermula di $Q$
yang menyentuh sebuah sumbu akan melalui sebuah titik salah satu
\omterm{def:b3:groups:action}{orbit} itu lalu, menurut \cref{cor:b3:ode:uniqueness}, berimpit
dengannya --- yang mustahil, sebab yang satu hidup pada sumbunya sedangkan yang lain
tidak. Jadi $Q$ invarian dalam kedua arah waktunya.
Untuk kesetimbangan di $Q$: $x > 0$ memaksa $\alpha - \beta y = 0$ dan
$y > 0$ memaksa $\delta x - \gamma = 0$: sehingga hanya satu titik
$(x_*, y_*) = \bigl(\frac\gamma\delta,
\frac\alpha\beta\bigr)$.

\textbf{18.} Sepanjang sebuah solusi,
\[
\dot H = \Bigl(\delta - \frac\gamma x\Bigr)x' +
\Bigl(\beta - \frac\alpha y\Bigr)y'
= (\delta x - \gamma)(\alpha - \beta y) +
(\beta y - \alpha)(\delta x - \gamma) = 0 .
\]
Sedangkan $f(x) = \delta x - \gamma\ln x$ mempunyai $f'' = \gamma/x^2 > 0$,
dengan $f'$ yang lenyap hanya di $x_*$, dan $f \to +\infty$ baik di
$0^+$ maupun di $+\infty$: jadi cembung sejati dan meledak, berminimum
$f(x_*)$; dan demikian pula $g(y) = \beta y - \alpha\ln y$, berminimum
$g(y_*)$. Sehingga $H \geq h_*$ dengan kesamaannya hanya di $(x_*,
y_*)$, dan setiap subaras $\{H \leq h\}\cap Q$ bersifat \omterm{def:b3:topology:compact}{kompak}:
sebab $f(x) \leq h - g(y_*)$ mengurung $x$ ke sebuah selang \omterm{def:b3:topology:compact}{kompak}
$\intoo0{+\infty}$ menurut keledakannya, demikian pula $y$, dan himpunannya
tertutup di $\R^2$ karena $H \to +\infty$ di batas
$Q$. Lalu sebuah \omterm{thm:b3:ode:maximal}{solusi maksimal} tinggal pada himpunan arasnya yang \omterm{def:b3:topology:compact}{kompak}, sehingga
ia tak dapat meninggalkan setiap \omterm{def:b3:topology:compact}{kompak} dalam waktu berhingga:
jadi \cref{thm:b3:ode:maximal} membuatnya global.

\textbf{19.} Tetapkan $h > h_*$ lalu tetapkan $c = h - g(y_*) >
f(x_*)$. Lalu karena $f$ turun sejati dari $+\infty$ ke
$f(x_*)$ pada $\intoc0{x_*}$ dan naik sejati kembali ke
$+\infty$ pada $\intco{x_*}{+\infty}$, persamaan $f(x) = c$
berakar tepat dua kali $x_- < x_* < x_+$, dan $\{f \leq c\} =
\intcc{x_-}{x_+}$. Lalu untuk $x \in \intoo{x_-}{x_+}$: $g(y) = h -
f(x) > g(y_*)$ berakar tepat dua kali $y_-(x) < y_* <
y_+(x)$, yang \omterm{def:b3:topology:continuity}{kontinu} terhadap $x$ (sebab balikan pembatasan $g$ yang
monoton sejati dan \omterm{def:b3:topology:continuity}{kontinu} pada kedua sisi
$y_*$), dengan $y_\pm(x) \to y_*$ saat $x \to x_\pm$; sedangkan di $x =
x_\pm$ satu-satunya solusinya adalah $y = y_*$. Jadi $\{H = h\}\cap
Q$ berupa gabungan grafik $y_+$ dan $y_-$ atas
$\intcc{x_-}{x_+}$, yang direkatkan di $(x_\pm, y_*)$: yakni sebuah kurva tertutup
di sekitar $(x_*, y_*)$ --- yaitu analogi oval
bandulnya.

\textbf{20.} Misalkan $C_h = \{H = h\}\cap Q$ dengan $h > h_*$: maka
satu-satunya kesetimbangan $Q$ berada di luar $C_h$, sehingga medannya tak pernah
lenyap padanya. Untuk tandanya: $x' = \beta x(y_* - y)$ dan $y' =
\delta y(x - x_*)$: sehingga geraknya ke kanan di bawah garis $y
= y_*$, ke atas di kanan $x = x_*$, ke kiri di atasnya, dan ke bawah di
kirinya --- yakni peredaran berlawanan arah jarum jam. Lalu ikutilah
solusinya dari sebuah titik $(x_0, y_-(x_0))$ pada cabang bawah
terbukanya, yang $x' > 0$ padanya: maka waktu untuk mencapai sudut kanannya
$B = (x_+, y_*)$ adalah
\[
\int_{x_0}^{x_+}\frac{\dd x}{\beta x\,\bigl(y_* -
y_-(x)\bigr)} .
\]
Lalu di dekat $x_+$, pilihlah $x' \in \intoo{x_*}{x_+}$; sehingga untuk $x \in
\intcc{x'}{x_+}$, $f(x_+) - f(x) \geq f'(x')\,(x_+ - x)$
(sebab $f'$ naik dan positif setelah $x_*$), sedangkan
hubungan arasnya dan ketaksamaan Taylor memberikan $g(y_-(x)) -
g(y_*) \leq \frac12\,\bigl(\max g''\bigr)\,(y_* -
y_-(x))^2$ pada rentang-$y$ \omterm{def:b3:topology:compact}{kompak} $C_h$: sehingga $y_* -
y_-(x) \geq c\,\sqrt{x_+ - x}$ dengan $c > 0$, dan
integrannya $O\bigl((x_+ - x)^{-1/2}\bigr)$: jadi \omterm{def:b3:lebesgue:l1}{terintegralkan} ---
yakni kekonvergenan pertanyaan 5, yang dialihkan. Sedangkan di tempat lain pada
cabangnya integrannya \omterm{def:b3:topology:continuity}{kontinu}. Jadi $B$ tercapai dalam
waktu berhingga; lalu di sana $y' = \delta y_*(x_+ - x_*) > 0$, sehingga
\omterm{def:b3:groups:action}{orbitnya} memasuki daerah $x > x_*$, $y > y_*$, mendaki ke
sudut atasnya $(x_*, y_+)$ lewat taksiran setangkupnya (dengan peran
$f$ dan $g$ yang ditukar), dan seterusnya mengelilingi keempat busurnya:
sehingga setelah waktu berhingga $T > 0$ solusinya kembali ke
titik awalnya. Lalu menurut \cref{cor:b3:ode:uniqueness} ia
berperiode $T$ --- yakni argumen pertanyaan 4, kata demi kata.

\textbf{21.} Pada sebuah \omterm{def:b3:groups:action}{orbit} berperiode $T$ di $Q$, $t \mapsto \ln
x(t)$ bersifat $\mathcal C^1$ dan berperiode $T$, sehingga
\[
0 = \int_0^T(\ln x)'\,\dd t
= \int_0^T(\alpha - \beta y)\,\dd t
= \alpha T - \beta\int_0^Ty\,\dd t ,
\]
yang memberikan $\frac1T\int_0^Ty = \frac\alpha\beta$; dan demikian pula $0 =
\int_0^T(\ln y)' = \delta\int_0^Tx\,\dd t - \gamma T$ memberikan
$\frac1T\int_0^Tx = \frac\gamma\delta$. Jadi rata-rata waktunya
adalah nilai kesetimbangannya, bagi setiap \omterm{def:b3:groups:action}{orbitnya} tanpa memandang
amplitudonya --- yakni hukum kekekalan yang tak seorang pun memasukkannya dengan tangan.

\textbf{22.} Dengan pemanenannya sistemnya kembali berbentuk
Lotka--Volterra, dengan parameter $\alpha -
\varepsilon$, $\beta$, $\gamma + \varepsilon$, dan $\delta$ (sebab
kesetimbangan \omterm{def:b3:topology:interior}{interiornya} bertahan karena $\varepsilon <
\alpha$). Lalu pertanyaan 21 yang diterapkan pada sistem barunya:
\[
\bar x = \frac{\gamma + \varepsilon}{\delta}
\ \ (\text{rata-rata mangsanya naik}),
\qquad
\bar y = \frac{\alpha - \varepsilon}{\beta}
\ \ (\text{rata-rata pemangsanya turun}) :
\]
sehingga pemanenan yang tak pandang bulu menggeser keseimbangannya ke arah
mangsanya. Sedangkan data d'Ancona membacanya secara terbalik: sebab perangnya memotong
penangkapan ikannya, $\varepsilon$ turun, sehingga rata-rata pemangsanya
$(\alpha - \varepsilon)/\beta$ naik dan rata-rata mangsanya
turun --- yakni pecahan hiu yang lebih besar pada tangkapannya, tepat seperti yang
dicatat pasar ikan Adriatik. Inilah asas Volterra,
yakni mekanisme yang sama di balik paradoks pestisidanya:
sebab memangkas kedua aras trofiknya menguntungkan aras yang dimakan.

\textbf{23.} Resepnya, pada kedua kalinya: (i) sebuah integral pertama
--- yakni $E$ bagi bandulnya dan $H$ di sini, yang ditemukan dengan memisahkan
$\dd y/\dd x$ --- meruntuhkan bidangnya menjadi kurva; (ii)
keledakan dan kekompakan himpunan arasnya memberikan keberadaan
global lewat \cref{thm:b3:ode:maximal}; (iii)
geometri arasnya --- yakni oval, dari bentuk $\cos$
di sana dan kecembungan sejati $f$ dan $g$ di sini --- terbaca
dari integralnya, bukan dari \omterm{ex:b3:ode:planeclassification}{alirannya}; (iv) sebuah medan tak lenyap
pada oval \omterm{def:b3:topology:compact}{kompaknya} ditambah kesingularan sudut yang \omterm{def:b3:lebesgue:l1}{terintegralkan}
memaksa waktu kembali yang berhingga; (v) ketunggalannya
(\cref{cor:b3:ode:uniqueness}) mengubah kembalinya menjadi
keperiodikan dan melarang kedatangan pada kesetimbangannya dalam waktu berhingga;
(vi) sedangkan buah hasilnya --- yakni uraian periode dan hukum rata-rata
--- datang dari teorema kekonvergenan yang diterapkan pada
integral yang dihasilkannya. Jadi baik $x'' = -\sin x$ maupun
Lotka--Volterra tak memiliki solusi bentuk tertutup yang dasar
(sebab integral eliptik pada yang satu, dan kurva aras transenden
pada yang lain), dan tak pernah satu pun diperlukan: sebab persamaannya
sendiri, yang diinterogasi secara kualitatif, menyerahkan seluruh
geraknya.

\textbf{24.} Pada sebuah aras energi $E_0 > 1$, $y^2 = 2(E_0 +
\cos x) > 0$: sehingga ekstremum $y^2$-nya di $\cos x = \pm1$, yang memberikan
batas yang dinyatakan itu, dan tercapai di $x \equiv 0, \pi$. Untuk rata-rata
waktunya atas satu periode $\tau = \tau(E_0)$: $x$ maju sejauh
tepat $2\pi$, sehingga
\[
\frac1\tau\int_0^\tau y\,\dd t = \frac{x(\tau) -
x(0)}{\tau} = \frac{2\pi}\tau .
\]
Sedangkan nisbah laju ekstremnya adalah $\sqrt{\frac{E_0 + 1}{E_0 -
1}} = 1 + O(E_0^{-1}) \to 1$: sehingga pada energi tinggi riak
$\pm1$ potensialnya terabaikan terhadap $E_0$, dan bandulnya
berputar hampir seragam --- yakni papan cucinya yang mendatar.

\textbf{25.} Pada $E_0 \geq 1 + \delta$, integran
$\tau(E_0) = \int_{-\pi}^{\pi}\frac{\dd x}{\sqrt{2(E_0 +
\cos x)}}$ terdominasi oleh $(2\delta)^{-1/2}$ sedangkan
turunan-$E_0$-nya
\[
\partial_{E_0}\frac1{\sqrt{2(E_0 + \cos x)}}
= -\frac{1}{\bigl(2(E_0 + \cos x)\bigr)^{3/2}}
\]
oleh $(2\delta)^{-3/2}$, keduanya \omterm{def:b3:lebesgue:l1}{terintegralkan} pada
$\intcc{-\pi}\pi$ yang \omterm{def:b3:topology:compact}{kompak}: sehingga pendiferensialan di bawah integralnya
(\cref{thm:b3:lebesgue:paramdiff}) berlaku dan memberikan
$\tau'(E_0) < 0$ (sebab integrannya negatif sejati):
jadi turun sejati dan $\mathcal C^1$. Untuk limitnya: saat $E_0 \to
\infty$, $\tau \leq \frac{2\pi}{\sqrt{2(E_0 - 1)}} \to 0$;
sedangkan saat $E_0 \to 1^+$: tulislah $E_0 + \cos x = (E_0 - 1) +
2\cos^2\frac x2$; maka integrannya naik saat $E_0$
turun, sehingga lewat kekonvergenan monoton
\[
\tau(E_0) \;\nearrow\;
\int_{-\pi}^{\pi}\frac{\dd x}{2\,\abs{\cos\frac x2}}
= +\infty,
\]
dengan integral limitnya divergen di $x = \pm\pi$ (sebab di sana
$\abs{\cos\frac x2} \sim \frac{\abs{x \mp \pi}}2$, yakni
$\frac1{\abs\cdot}$ yang tak \omterm{def:b3:lebesgue:l1}{terintegralkan}): sehingga periodenya meledak
saat mendekati separatriksnya, yang cocok dengan $T(a) \to
\infty$ Bagian II dari sisi librasinya. Jadi sumbu
energinya terbaca: istirahat di $E_0 = -1$; librasi, $2\pi \nearrow
\infty$, pada $\intoo{-1}1$; separatriks yang tak berhingga lambat di
$E_0 = 1$; lalu rotasi, $\infty \searrow 0$, selebihnya. Yakni satu
integral, seluruh hidup bandulnya.

\end{solution}
